\documentclass[11pt]{article}
\usepackage{varwidth}
\usepackage{manfnt}
\usepackage[colorlinks=true,linkcolor=blue,citecolor=blue]{hyperref}
\usepackage[scaled]{beramono}
\usepackage{fullpage}
\usepackage[dvips]{graphicx} 
\usepackage{enumerate}
\usepackage{amsmath}
\usepackage{amsfonts}
\usepackage{amsthm}
\usepackage{amssymb}
\usepackage[usenames,dvipsnames]{xcolor}
\usepackage{multicol}
\usepackage{pifont}
\usepackage{tikz}
\usepackage{mathabx}
\usepackage{xparse}
\usepackage[scaled]{beramono}
\usepackage{ragged2e}
\usepackage{mdframed}
\usepackage{derivative}
\usepackage{xcolor}

\usepackage{algorithm}
\usepackage[noend]{algpseudocode}
\usepackage{minted}

\usetikzlibrary{arrows}


\renewcommand{\epsilon}{\varepsilon}

\newcommand{\ngroup}[1]{\left(#1 \right)}
\newcommand{\sgroup}[1]{\left[#1 \right]}
\newcommand{\cgroup}[1]{\left\{#1 \right\}}
\newcommand{\agroup}[1]{\left\langle #1 \right\rangle}


\newcommand{\bbN}{\mathbb{N}}

\newcommand{\set}[1]{\mathcal{#1}}
\newcommand{\setrange}[2]{\{#1, \ldots, #2\}}

\newcommand{\Prob}[1]{\Pr\left[#1 \right]}
\newcommand{\CondProb}[2]{ \Prob{#1 \,\big|\, #2}}
\newcommand{\abs}[1]{\left| #1 \right|}

\newcommand{\bA}{\mathbf{A}}
\newcommand{\bB}{\mathbf{B}}
\newcommand{\bC}{\mathbf{C}}
\newcommand{\bD}{\mathbf{D}}
\newcommand{\bE}{\mathbf{E}}
\newcommand{\bF}{\mathbf{F}}
\newcommand{\bG}{\mathbf{G}}
\newcommand{\bH}{\mathbf{H}}
\newcommand{\bI}{\mathbf{I}}
\newcommand{\bJ}{\mathbf{J}}
\newcommand{\bK}{\mathbf{K}}
\newcommand{\bL}{\mathbf{L}}
\newcommand{\bM}{\mathbf{M}}
\newcommand{\bN}{\mathbf{N}}
\newcommand{\bO}{\mathbf{O}}
\newcommand{\bP}{\mathbf{P}}
\newcommand{\bQ}{\mathbf{Q}} 
\newcommand{\bR}{\mathbf{R}}
\newcommand{\bS}{\mathbf{S}}
\newcommand{\bT}{\mathbf{T}}
\newcommand{\bU}{\mathbf{U}}
\newcommand{\bV}{\mathbf{V}}
\newcommand{\bW}{\mathbf{W}}
\newcommand{\bX}{\mathbf{X}}
\newcommand{\bY}{\mathbf{Y}}
\newcommand{\bZ}{\mathbf{Z}}

\newcommand{\Unif}{\mathrm{Unif}}

\newcommand{\bfa}{\mathbf{a}}
\newcommand{\bfb}{\mathbf{b}}
\newcommand{\bfc}{\mathbf{c}}
\newcommand{\bfd}{\mathbf{d}}
\newcommand{\bfe}{\mathbf{e}}
\newcommand{\bff}{\mathbf{f}}
\newcommand{\bfg}{\mathbf{g}}
\newcommand{\bfh}{\mathbf{h}}
\newcommand{\bfi}{\mathbf{i}}
\newcommand{\bfj}{\mathbf{j}}
\newcommand{\bfk}{\mathbf{k}}
\newcommand{\bfl}{\mathbf{l}}
\newcommand{\bfm}{\mathbf{m}}
\newcommand{\bfn}{\mathbf{n}}
\newcommand{\bfo}{\mathbf{o}}
\newcommand{\bfp}{\mathbf{p}}
\newcommand{\bfq}{\mathbf{q}}
\newcommand{\bfr}{\mathbf{r}}
\newcommand{\bfs}{\mathbf{s}}
\newcommand{\bft}{\mathbf{t}}
\newcommand{\bfu}{\mathbf{u}}
\newcommand{\bfv}{\mathbf{v}}
\newcommand{\bfw}{\mathbf{w}}
\newcommand{\bfx}{\mathbf{x}}
\newcommand{\bfy}{\mathbf{y}}
\newcommand{\bfz}{\mathbf{z}}
\newcommand{\expect}[1]{\mathbb{E}\left[#1\right]}
\newcommand{\Var}[1]{\mathrm{Var}\left(#1\right)}
\newcommand{\Exp}[1]{\expect{#1}}
\newcommand{\E}{\mathbb{E}}
\newcommand{\V}{\mathrm{Var}}

\newcommand{\bin}[1]{\ensuremath{\{0,1\}^{#1}}}
\newcommand{\rsample}{\stackrel{\mbox{\tiny $\$$}}{\leftarrow}}
\newcommand{\setd}[2]{\left\{ #1 \,:\, #2 \right\}}


\renewcommand{\labelitemi}{-}
\newcommand{\bits}{\{0,1\}}

\renewcommand{\Pr}{\mathbb{P}}
\renewcommand{\P}[1]{\Prob{#1}}

\newcommand{\bbR}{\mathbb{R}}

\theoremstyle{plain}
\newtheorem{thm}{Theorem}[section]
\newtheorem{lemma}[thm]{Lemma}
\newtheorem{proposition}[thm]{Proposition}
\newtheorem*{corollary}{Corollary}
\theoremstyle{definition}
\newtheorem{defn}{Definition}[section]

\newcommand{\Adv}{\mathrm{Adv}}


\newcommand{\IBE}{\mathsf{IBE}}
\newcommand{\Extract}{\mathsf{Extract}}
\newcommand{\pp}{\mathsf{pp}}
\newcommand{\msk}{\mathsf{msk}}
\newcommand{\ID}{\mathrm{ID}}
\newcommand{\id}{\mathrm{id}}
\newcommand{\IDspace}{\ensuremath{\mathcal{ID}}}

\newcommand{\IBEEnc}{\mathsf{IBEEnc}}
\newcommand{\IBEGen}{\mathsf{IBEGen}}
\newcommand{\IBEDec}{\mathsf{IBEDec}}
\newcommand{\IBEExtract}{\mathsf{IBEExtract}}
\renewcommand{\SS}{\ensuremath{\mathsf{SS}}}


\newcommand{\PKE}{\mathsf{PKE}}
\newcommand{\Gen}{\mathsf{Gen}}
\newcommand{\Enc}{\mathsf{Enc}}
\newcommand{\Dec}{\mathsf{Dec}}
\newcommand{\pk}{\mathsf{pk}}
\newcommand{\sk}{\mathsf{sk}}
\newcommand{\getsr}{\rsample}
\newcommand{\advA}{\mathcal{A}}
\newcommand{\indcpa}{\mathrm{ind}\mbox{-}\mathrm{cpa}}
\newcommand{\true}{\texttt{true}}
\newcommand{\false}{\texttt{false}}

\renewcommand{\Pr}{\mathrm{Pr}}
\newcommand{\prf}{\mathsf{prf}}
\newcommand{\dist}{\mathsf{dist}}
\newcommand{\pred}{\mathrm{pred}}
\newcommand{\bbG}{\mathbb{G}}
\renewcommand{\mod}{\textrm{ mod }}


\begin{document}


\title{{\bf Homework 1 Solutions} }
\author{Names: TODO}
\date{}

\maketitle

\section*{Task 1 -- Negligible Functions}

\noindent Recall that a function $\epsilon: \bbN \to \bbR_{\ge 0}$ is {\em
  negligible} if for all $c \in \bbN$, there exists $\lambda_0(c)$ such that
$\epsilon(\lambda) \leq \lambda^{-c}$ for all $\lambda \ge \lambda_0(c)$.

\begin{enumerate}
\item Which ones of the following functions are negligible, and which
  ones are not? (Logarithms here are base two.) Prove your claim! (Some of these functions have been
  claimed to be negligible in class -- here, you want to give a proof.)
\begin{multicols}{3}
\begin{itemize}
  \item $\epsilon_1(\lambda) = 2^{-\lambda}$
  \item $\epsilon_2(\lambda) = 2^{-\log^2(\lambda)}$
  \item $\epsilon_3(\lambda) = 1/2 + 1/2(-1)^\lambda$
\end{itemize}
\end{multicols}

\item  Show that for any two negligible functions $\epsilon_1, \epsilon_2$,
  and any polynomial $p$ which is non-negative (i.e., $p(\lambda) \geq 0$ on
  all $\lambda \in \bbN$) the following functions are negligible:
  \begin{multicols}{3}
    \begin{itemize}
      \item $\lambda \mapsto \epsilon_1(\lambda) + \epsilon_2(\lambda)$
      \item $\lambda \mapsto \epsilon_1(\lambda) \cdot \epsilon_2(\lambda)$
      \item $\lambda \mapsto p(\lambda) \cdot \epsilon_1(\lambda)$
    \end{itemize}
   \end{multicols}
   \noindent {\bf Hint:} Note that $a \cdot \lambda^c < \lambda^{c+1}$ for any $a
   \in \mathbb{R}$ and $c \in \bbN$, and large enough $\lambda$. Recall that
   a polynomial (with degree $d$) is a function of form $p(\lambda) = a_0 +
   a_1 \lambda + a_2 \lambda^2 + \cdots + a_d \lambda^d$.
 \end{enumerate}

\newpage

\section*{Task 2 -- Distinguishing Advantages}

\noindent We have seen two approaches to defining PRF
security of a keyed function $F: \bits^\lambda \times
\bits^{n(\lambda)} \to \bits^{\ell(\lambda)}$.
\begin{itemize}
\item First, via the game $\mathsf{PRF}_{F}(\advA,\lambda)$, and then defining
  the advantage as $$\Adv^{\mathsf{prf}}_F(\advA,\lambda) = \abs{2 \cdot
    \Prob{\mathsf{PRF}_{F}(\advA,\lambda) \Rightarrow \true} - 1}\;.$$
\item Second, via the two games  $\mathsf{PRF}_{F,0}(\advA,\lambda)$ and
  $\mathsf{PRF}_{F,1}(\advA,\lambda)$, and then defining the advantage as
  $$\Adv^{\mathsf{prf}'}_F(\advA,\lambda) = \abs{
    \Prob{\mathsf{PRF}_{F,0}(\advA,\lambda) \Rightarrow \true}-\Prob{\mathsf{PRF}_{F,1}(\advA,\lambda) \Rightarrow \true} }\;.$$
\end{itemize}
Prove that the following identity holds for all
   distinguishers $\advA$:
   \begin{displaymath}
     \Adv^{\mathsf{prf}}_F(\advA,\lambda) =
     \Adv^{\mathsf{prf}'}_F(\advA,\lambda) \;.
   \end{displaymath}
   {\bf Hint.} First argue that
   \begin{displaymath}
\Prob{\mathsf{PRF}_{F}(\advA,\lambda) \Rightarrow \true} = \frac{1}{2}
\Prob{\mathsf{PRF}_{F,0}(\advA,\lambda) \Rightarrow \false} +
\frac{1}{2} \Prob{\mathsf{PRF}_{F,1}(\advA,\lambda) \Rightarrow \true} \;.
   \end{displaymath}

\newpage


\section*{Task 3 -- Message Authentication Codes}

 \noindent Let $F: \bits^\lambda \times \bits^\lambda \to \bits^\lambda$ be a
 UF-CMA secure MAC. We now consider a few keyed functions constructed using $F$ in the following questions. If the construction is UF-CMA secure, please give a proof. If not, please show an attack.
 \begin{enumerate}
 \item Consider $F_1: \bits^\lambda \times
   \bits^{\lambda} \to \bits^{\lambda + 1}$ such that
   \begin{displaymath}
     F_1(K, X) = (F(K, X) , 1) \;,
   \end{displaymath}
   for all $K \in \bits^{\lambda}$ and $X \in \bits^\lambda$. Is $F_1$ UF-CMA secure?
 \item Consider $F_2: \bits^{2\lambda} \times
   \bits^{\lambda} \to \bits^{2\lambda}$ such that
   \begin{displaymath}
     F_2((K_1,K_2), X) = (F(K_1, X), F(K_2, X))\;,
   \end{displaymath}
   for all $K_1,K_2 \in \bits^\lambda$ and $X \in \bits^\lambda$. Is $F_2$ UF-CMA secure?
  \item Consider $F_3: \bits^\lambda \times
   \bits^{2\lambda} \to \bits^{2\lambda}$ such that
   \begin{displaymath}
     F_3(K, (X_1, X_2)) = (F(K, X_1), F(K, X_2)) \;,
   \end{displaymath}
   for all $K \in \bits^\lambda$ and $X_1, X_2 \in \bits^\lambda$. Is $F_3$ UF-CMA secure?
 \end{enumerate}
Note that here we use vector notation $(X,Y)$ to denote the
concatenation of two strings $X$ and $Y$.

\newpage

\section*{Task 4 -- Pseudorandom Functions}

 \noindent Let $F: \bits^\lambda \times \bits^\lambda \to
 \bits^\lambda$ be a keyed function which is a
 PRF. We now consider a few keyed functions constructed using $F$ in
 the following questions. If the resulting construction is a PRF, please give a proof. If not, please show an attack.
 \begin{enumerate}
 \item  Consider $F_1: \bits^\lambda \times
   \bits^{2\lambda} \to \bits^{\lambda}$ such that
   \begin{displaymath}
     F_1(K, (X_1, X_2)) = F(K, X_1) \oplus F(K, X_2) \;,
   \end{displaymath}
   for all $K \in \bits^{\lambda}$ and $X_1, X_2 \in \bits^\lambda$. Is $F_1$ a PRF?
 \item Consider $F_2: \bits^{2\lambda} \times
   \bits^{2\lambda} \to \bits^{\lambda}$ such that
   \begin{displaymath}
     F_2((K_1, K_2), (X_1, X_2)) = F(K_1, X_1) \oplus F(K_2, X_2) \;,
   \end{displaymath}
   for all $K_1,K_2 \in \bits^\lambda$ and $X_1, X_2 \in \bits^\lambda$. Is $F_2$ a
   PRF?
  \item Consider $F_3: \bits^\lambda \times
   \bits^\lambda \to \bits^\lambda$ such that
   \begin{displaymath}
     F_3(K, X) = F(K, X) \oplus X \;,
   \end{displaymath}
   for all $K \in \bits^\lambda$ and $X \in \bits^\lambda$. Is $F_3$ a PRF?
   \item Consider $F_4: \bits^\lambda \times
   \bits^{\lambda-1} \to \bits^{2\lambda}$ such that
   \begin{displaymath}
     F_4(K, X) = (F(K, (X,0)), F(K, (X, 1))) \;,
   \end{displaymath}
   for all $K \in \bits^\lambda$ and $X \in \bits^{\lambda - 1}$. Is $F_4$ a PRF?
 \end{enumerate}
Note that here we use vector notation $(X,Y)$ to denote the
concatenation of two strings $X$ and $Y$.

\newpage


\newcommand{\advG}{\mathcal{G}}
\newcommand{\Eval}{\mathbf{E}}
\newcommand{\ks}{\mathcal{K}}
\newcommand{\ms}{\mathcal{X}}
\newcommand{\os}{\mathcal{Y}}
\newcommand{\leftkeys}{\mathsf{Keys}}
\newcommand{\farkeys}{\mathsf{FarKeys}}
\newcommand{\diff}{\mathsf{Diff}}

\section*{Breaking UF-CMA
    Security in Exponential Time}

 \noindent Let $F:\bits^{\lambda} \times \bits^{n(\lambda)} \to \bits^{\ell(\lambda)}$
 be a keyed function. Consider the following adversary $\advG$ against UF-CMA security of $F$.
 \begin{center}
   \begin{minipage}[c]{0.5\textwidth}
       \underline{{\bf Adversary} $\advG^\Eval(1^{\lambda})$:} \\[0.2ex]
       $X_1,\ldots,X_q \getsr \bits^{n(\lambda)}$\\
       For all $i\in [q]$ do $T_i\gets \Eval(X_i)$\\
       $\leftkeys\gets \{K' \in \bits^{\lambda}|~\forall i\in [q]: F(K',X_i)=T_i\}$\\
       $K'\getsr \leftkeys$\\
       $X^*\getsr \bits^{n(\lambda)} \setminus \{X_1,\ldots,X_q\}$\\
       $T^*\gets F(K',X^*)$\\
       {\bf return} $(X^*,T^*)$
   \end{minipage}

 \end{center}
Consider now an execution of game
$\mathsf{UF}\mbox{-}\mathsf{CMA}_{F}(\mathcal{G}, \lambda)$. Let $K$
be the true key that is sampled by the game. Let $\farkeys_K$ be the
set of keys $K'$ such that $F(K', \cdot)$ and $F(K,\cdot)$ differ on at least
half fraction of the inputs. Formally, we define the quantity
\begin{align*}
  \farkeys_K =\{K' \in \bits^{\lambda}|~ |\diff_{K, K'}| \geq 2^{n(\lambda) - 1}\}\;,
\end{align*}
where $\diff_{K, K'} := \{X\in \bits^{n(\lambda)}  |~ F(K, X) \neq F(K', X)\}$.

    \begin{enumerate}

 \item Prove that
 \begin{align*}
   \Prob{\leftkeys \cap \farkeys_K \neq \emptyset} \leq  \frac{2^\lambda}{2^q}\;.
 \end{align*}
 \item Prove that conditioned on $\leftkeys \cap \farkeys_K
   =\emptyset$, $F(K, X^*) = F(K', X^*)$ is going to hold with high
   probability for $K' \getsr \leftkeys$, i.e.,
 \begin{align*}
  \Prob{F(K, X^*) = F(K', X^*) \;|\; \leftkeys \cap \farkeys_K
   =\emptyset} \geq \frac{2^{n(\lambda) - 1}-q}{2^{n(\lambda)}}\;.
\end{align*}
 \item Prove that for $\lambda +1\leq q\leq 2^{n(\lambda) - 2}$,
 \begin{align*}
   \Adv^{\mathsf{uf}\mbox{-}\mathsf{cma}}_{F}(\advG, \lambda)\geq \frac{1}{8}\;.
 \end{align*}
 \end{enumerate}

\end{document}
